\documentclass[11pt,a4paper]{article}
\usepackage[utf8]{inputenc}
\usepackage[T1]{fontenc}
\usepackage[french]{babel}
\usepackage{lmodern}
\usepackage{amsmath,amssymb}
\usepackage[margin=2.5cm]{geometry}
\usepackage{microtype}
\usepackage{hyperref}
\usepackage{booktabs}
\usepackage{listings}
\lstset{basicstyle=\ttfamily\small, frame=single, breaklines=true,
  columns=fullflexible, keepspaces=true}
\title{Projet MLSF : résolution du jeu Meowdoku\\
  \large Question 4 : heuristiques de branchement (DPLL$_H$)}
\author{}
\date{}

\begin{document}
\maketitle

\section{Objectif et organisation}

L'algorithme 3 de l'énoncé est l'algorithme 2 dans lequel la fonction
\texttt{RandomChoice} (ligne 17) est remplacée par \texttt{Heuristics(F, H)}.
La question demande d'implémenter au moins quatre heuristiques $H$, dont au
moins une fondée sur le choix d'un littéral, une qui tient compte de la taille
des clauses et une dynamique, puis de vérifier qu'elles donnent les mêmes
résultats que l'algorithme 2.

Nous avons implémenté \textbf{toutes} les heuristiques de l'énoncé, dans
leurs deux variantes statique et dynamique, ainsi qu'une heuristique de notre
choix (H8), soit 16 heuristiques :
\begin{center}
\begin{tabular}{lll}
  \toprule
  Nom & Heuristique & Sélection \\
  \midrule
  \texttt{H0} & CA : choix aléatoire (algorithme 2) & variable \\
  \texttt{H1} & PVL : première variable libre, valeur 1 & variable \\
  \texttt{H2s}, \texttt{H2d} & VPPM : variable la plus fréquente & variable \\
  \texttt{H3s}, \texttt{H3d} & LP : littéral le plus fréquent & littéral \\
  \texttt{H4s}, \texttt{H4d} & VFCM (MOMS) : clauses de taille minimale & variable \\
  \texttt{H5s}, \texttt{H5d} & LFCM : clauses de taille minimale & littéral \\
  \texttt{H6s}, \texttt{H6d} & VJW : poids de Jeroslow-Wang & variable \\
  \texttt{H7s}, \texttt{H7d} & LJW : poids de Jeroslow-Wang & littéral \\
  \texttt{H8s}, \texttt{H8d} & JWP : Jeroslow-Wang, littéraux positifs
    (section~\ref{sec:h8}) & littéral \\
  \bottomrule
\end{tabular}
\end{center}
Le suffixe \texttt{s} désigne la variante statique, \texttt{d} la variante
dynamique. Les contraintes de l'énoncé sont donc couvertes plusieurs fois :
choix d'un littéral (H3, H5, H7, H8), taille des clauses (H4 à H8),
variantes dynamiques (H2d à H8d).

Le code se répartit ainsi :
\begin{itemize}
  \item \texttt{heuristics.py} (nouveau) : les heuristiques et la fonction
    \texttt{dpll\_h(clauses, heuristic)} qui implémente l'algorithme 3 ;
  \item \texttt{dpll.py} : la fonction de branchement de \texttt{dpll}
    devient un paramètre, et un compteur du temps passé à choisir est ajouté
    (préparation de la question 5) ;
  \item \texttt{tests/test\_heuristics.py} (nouveau) : tests unitaires ;
  \item \texttt{validate.py} : option \texttt{--heuristics} pour comparer
    les heuristiques à l'algorithme 2 sur des fichiers DIMACS.
\end{itemize}

\section{Intégration dans DPLL : une seule implémentation}
\label{sec:integration}

\paragraph{Choix.} Plutôt que d'écrire une seconde fonction DPLL, la
fonction \texttt{dpll} de la question 2 reçoit un paramètre facultatif
\texttt{choose}, une fonction qui prend la formule courante et renvoie un
littéral :
\begin{lstlisting}[language=Python]
def dpll(clauses, rng=None, stats=None, choose=None):
    if choose is None:
        rng = rng if rng is not None else random.Random()
        choose = lambda current: random_choice(current, rng)
    ...
        literal = choose(clauses)      # ligne 17 de l'algorithme 3
\end{lstlisting}
Comme à la question 2, le couple $(p, b)$ est codé par un seul littéral :
$p$ si $b = 1$, $-p$ si $b = 0$ ; la branche \og sinon \fg{} utilise son
opposé. \texttt{heuristics.dpll\_h} construit la fonction \texttt{choose}
correspondant au nom de l'heuristique et appelle \texttt{dpll}.

\paragraph{Justification.} L'algorithme 3 ne diffère de l'algorithme 2 que
par la ligne 17 : partager le code garantit que toutes les autres étapes
(clause vide, clause unitaire, littéral pur, construction de la valuation)
sont exactement celles qui ont été validées à la question 2. Une copie de
\texttt{dpll} aurait pu diverger lors d'une correction future. Sans
paramètre \texttt{choose}, le comportement de la question 2 est inchangé ;
un test vérifie qu'avec la même graine, \texttt{H0} donne exactement le même
résultat et le même nombre de décisions que l'algorithme 2.

\paragraph{Fonctions plutôt que classes.} Une heuristique est une simple
fonction \texttt{choose(clauses)}. Les variantes statiques ont besoin d'un
état (l'ordre calculé au départ) : il est capturé par une fermeture
(fonction définie à l'intérieur de \texttt{make\_static}). Cela évite de
définir une hiérarchie de classes pour un comportement d'une ligne.

\section{Ce que voit l'heuristique}

\paragraph{La formule n'a que des clauses de taille au moins 2.} La ligne 17
n'est atteinte que si la formule courante ne contient ni clause vide, ni
clause unitaire, ni littéral pur, et n'est pas vide. Elle contient donc au
moins une variable, et $k = \min_{C \in F} |C| \geq 2$. Les heuristiques
n'ont pas à traiter de cas particulier.

\paragraph{\og Variable non affectée \fg{} = variable présente dans la
formule courante.} Notre DPLL ne stocke pas d'affectation pendant la
descente : il simplifie la formule (\texttt{simplify}). Une variable absente
de la formule courante est soit déjà affectée, soit apparue seulement dans
des clauses déjà satisfaites ; dans les deux cas la brancher serait inutile
(la formule ne changerait pas). Les heuristiques choisissent donc parmi les
variables de la formule courante, comme \texttt{random\_choice} à la
question 2. C'est aussi ce qui rend H1 (PVL) correcte : c'est la plus petite
variable de la formule courante.

\section{Une structure commune à H2--H8}

Toutes les heuristiques H2 à H8 comptent des occurrences de littéraux, chaque
occurrence étant pondérée par un poids qui ne dépend que de la taille de sa
clause. Elles ne diffèrent que par ce poids et par le mode de sélection :
\begin{center}
\begin{tabular}{lll}
  \toprule
  Heuristiques & Poids d'une clause $C$ & Fonction \\
  \midrule
  H2, H3 & $1$ & \texttt{weight\_one} \\
  H4, H5 & $1$ si $|C| = k$, $0$ sinon & \texttt{weight\_min\_size} \\
  H6, H7, H8 & $2^{-|C|}$ & \texttt{weight\_jeroslow\_wang} \\
  \bottomrule
\end{tabular}
\end{center}
Une seule fonction calcule ces scores :
\begin{lstlisting}[language=Python]
def literal_scores(clauses, weight):
    min_size = min((len(clause) for clause in clauses), default=0)
    scores = Counter()
    for clause in clauses:
        w = weight(len(clause), min_size)
        if w:
            for literal in clause:
                scores[literal] += w
    return scores
\end{lstlisting}
Avec le poids 1, \texttt{scores[l]} vaut $occ(l, F)$, et le score d'une
variable est $\texttt{scores[p]} + \texttt{scores[-p]} = \#(p, F)$. Avec le
poids de H4, on obtient $occ_k(l, F)$ et $\#_k(p, F)$ ; avec $2^{-|C|}$, les
sommes de Jeroslow-Wang $\sum_k occ_k(l, F)\,2^{-k}$ et
$\sum_k \#_k(p, F)\,2^{-k}$.

\paragraph{Deux modes de sélection, décrits par une clé de tri.} Le littéral
choisi est celui qui maximise une clé :
\begin{itemize}
  \item \textbf{variable} (H2, H4, H6) : la clé de $\pm p$ est
    $(\#(p),\ -p,\ \text{$\pm p$ est la polarité majoritaire})$ ; on obtient
    la variable de plus grand score avec sa polarité majoritaire, $1$ si
    $\#^+ \geq \#^-$ comme dans l'énoncé ;
  \item \textbf{littéral} (H3, H5, H7) : la clé de $l$ est
    $(\texttt{scores[l]},\ -|l|,\ l > 0)$.
\end{itemize}
\paragraph{Égalités.} En cas d'égalité de score, la plus petite variable
l'emporte, puis la polarité positive (le $-p$ et le $l>0$ des clés). Le choix
est donc déterministe : une exécution ne dépend pas de l'ordre interne des
\texttt{set} Python et peut être reproduite exactement, ce qui facilite le
débogage et les mesures. Seule H0 utilise le hasard (avec une graine).

\paragraph{Justification.} Factoriser les six heuristiques de l'énoncé en
\og poids + sélection \fg{} tient en 200 lignes commentaires compris, met en
évidence ce qui les distingue réellement et garantit qu'un même calcul de
scores est utilisé partout (un bug ne peut pas toucher une seule heuristique
sans être visible sur les autres). Ajouter une heuristique (H8) a demandé
trois lignes.

\section{Variantes statiques et dynamiques}

\paragraph{Dynamique.} Les scores sont recalculés sur la formule courante à
chaque décision, ce qui coûte un parcours de la formule ($O(|F|)$, le même
ordre que \texttt{simplify}) :
\begin{lstlisting}[language=Python]
def choose(clauses):
    key = make_key(literal_scores(clauses, weight))
    return max(literals_of(variables_of(clauses)), key=key)
\end{lstlisting}

\paragraph{Statique.} Comme le suggère l'énoncé, un prétraitement calcule
les scores une seule fois sur la formule initiale et trie les littéraux par
clé décroissante ; on ne garde que le premier littéral de chaque variable.
Choisir revient ensuite à prendre le premier littéral de cet ordre dont la
variable est encore présente dans la formule courante. Sur la formule
initiale, la variante statique fait donc exactement le même choix que la
variante dynamique (c'est testé) ; elles divergent ensuite, quand la formule
a été simplifiée.

Le choix statique parcourt encore la formule pour connaître ses variables :
nous l'avons accepté pour garder la représentation de la question 2 (une
liste de clauses) sans structure supplémentaire. Ce parcours est bien moins
coûteux que le calcul des scores, et la question 6 pourra le supprimer avec
une affectation explicite.

\paragraph{H4 statique.} Pour la variante statique, $k$ est la taille
minimale des clauses de la formule initiale ; les variables qui n'apparaissent
dans aucune clause de taille $k$ ont un score nul et sont rangées après les
autres, par numéro croissant.

\section{Vérification sur l'exemple de l'énoncé}

Pour $F = (x_1 \lor x_2 \lor \lnot x_4) \land (x_2 \lor x_3) \land
(x_1 \lor \lnot x_2 \lor x_4 \lor \lnot x_5)$, les scores calculés à la main
sont :
\begin{center}
\begin{tabular}{lccccc}
  \toprule
  & $x_1$ & $x_2$ & $x_3$ & $x_4$ & $x_5$ \\
  \midrule
  $\#^+ / \#^-$ & 2 / 0 & 2 / 1 & 1 / 0 & 1 / 1 & 0 / 1 \\
  $\#_2$ ($k = 2$) & 0 & 1 & 1 & 0 & 0 \\
  $\sum_k \#_k\,2^{-k}$ & $3/16$ & $7/16$ & $1/4$ & $3/16$ & $1/16$ \\
  \bottomrule
\end{tabular}
\end{center}
D'où les choix : H1 $\to x_1$ ; H2 $\to x_2 = 1$ ($\#(x_2) = 3$, maximum) ;
H3 $\to x_1$ ($occ(x_1) = occ(x_2) = 2$, égalité départagée par le numéro) ;
H4 et H5 $\to x_2 = 1$ ($x_2$ et $x_3$ à égalité dans la seule clause de
taille 2) ; H6 $\to x_2 = 1$ (poids $7/16$, la valeur donnée dans l'énoncé) ;
H7 et H8 $\to x_2$ (poids de $x_2$ : $1/8 + 1/4 = 3/8$). Ces choix sont
vérifiés par \texttt{test\_choix\_sur\_l\_exemple}.

\paragraph{Lecture de l'énoncé.} Quelques coquilles ont été interprétées :
\begin{itemize}
  \item pour H5, l'énoncé indique \og 1 si $l$ est de la forme $\lnot p$ \fg{}
    dans la première phrase, puis 0 dans la définition formelle : nous
    affectons la valeur qui rend $l$ vrai (0 pour $\lnot p$), seule lecture
    cohérente avec H3 et H7 ;
  \item pour H7, la formule écrit $\#_k(l, F)$ ; nous utilisons
    $occ_k(l, F)$, comme dans la phrase qui la précède ;
  \item dans l'exemple, la ligne \og $\#_2(x_2, F) = 1$, $\#^+(x_2, F) = 2$
    \fg{} mélange les notations ; nos valeurs de $\#^+_2$ et $\#^-_2$ sont
    1 et 0.
\end{itemize}

\section{Heuristique ajoutée : H8, Jeroslow-Wang positif}
\label{sec:h8}

\paragraph{Constat.} Les premières mesures (section~\ref{sec:mesures})
montrent que, sur les grilles Meowdoku, toutes les heuristiques H2 à H7 font
\emph{plus} de décisions que le choix aléatoire. L'explication vient de
l'encodage de la question 1 : les contraintes \og au plus un chat \fg{} (par
ligne, colonne, région et voisinage) donnent un très grand nombre de clauses
binaires $\lnot a \lor \lnot b$, toutes négatives ; seules les clauses
\og au moins un chat par région \fg{} contiennent des littéraux positifs, et
elles sont longues. Toute heuristique fondée sur la fréquence ou la
polarité majoritaire choisit donc presque toujours \og pas de chat dans
cette case \fg{}. Cette décision ne fait disparaître que quelques clauses et
ne déclenche presque aucune propagation. À l'inverse, placer un chat
(affecter une variable à 1) rend unitaires toutes les clauses binaires de la
case : la propagation élimine toute la ligne, la colonne, la région et les
voisins.

\paragraph{Définition.} H8 utilise les poids de Jeroslow-Wang (comme H7)
mais ne choisit que parmi les littéraux \emph{positifs} présents dans la
formule ; s'il n'en reste aucun, elle se comporte comme H7. Sa clé est
$(l > 0 \text{ et } occ(l) > 0,\ \texttt{scores[l]},\ -|l|,\ l > 0)$.
Sur Meowdoku, un littéral positif n'apparaît que dans la clause \og au moins
un chat \fg{} de sa région, de taille le nombre de cases encore libres de la
région ; son poids est $2^{-\text{taille}}$. H8 place donc un chat dans la
\textbf{région qui a le moins de cases possibles}. C'est le principe
classique \og variable la plus contrainte d'abord \fg{} des problèmes de
contraintes : si la région n'a que deux cases libres, l'une des deux
branches est la bonne et l'échec est détecté vite.

\paragraph{Pourquoi la garder générale.} H8 est définie pour toute formule
CNF, pas seulement pour Meowdoku : elle passe les mêmes tests que les autres
heuristiques et reste raisonnable sur les instances aléatoires (voir les
mesures). Elle est la candidate naturelle pour le solveur optimisé de la
question 6.

\section{Validation}
\label{sec:validation}

L'énoncé demande que les heuristiques donnent des résultats conformes à
l'algorithme 2 sur les données de test de la question 2. Trois niveaux de
vérification ont été mis en place.

\paragraph{Tests unitaires (\texttt{tests/test\_heuristics.py}).}
\begin{itemize}
  \item les comptages et le poids $7/16$ de l'exemple de l'énoncé, et le
    choix de chaque heuristique sur cet exemple ;
  \item la polarité minoritaire n'est jamais choisie par VPPM ; H8 préfère
    un littéral positif à un négatif de meilleur poids ;
  \item statique et dynamique coïncident sur la formule initiale et
    divergent après une simplification (exemple explicite) ;
  \item une variante statique ne renvoie jamais une variable sortie de la
    formule courante ;
  \item pour chacune des 16 heuristiques : cas limites (formule vide,
    clause vide, tautologie, exemple des lampes), 400 formules aléatoires
    comparées à une force brute qui essaie les $2^n$ valuations, grilles
    Meowdoku SAT de taille 4 à 8 et une grille UNSAT ; le modèle renvoyé
    est vérifié clause par clause ;
  \item \texttt{H0} reproduit exactement l'algorithme 2 à graine égale.
\end{itemize}
Tous les tests passent :
\begin{lstlisting}
uv run python -m unittest discover -s tests -t .
\end{lstlisting}

\paragraph{Instances SATLIB et instances publiques.} L'option
\texttt{--heuristics} de \texttt{validate.py} lance, en plus des
algorithmes 1 et 2, chaque heuristique demandée. Elle doit renvoyer le même
statut que l'algorithme 2 (et que le statut attendu) et, si SAT, un modèle
qui satisfait toutes les clauses. Une heuristique qui dépasse le délai n'est
pas une erreur mais elle est signalée.
\begin{lstlisting}
uv run python validate.py data/satlib --heuristics all --timeout 10
uv run python validate.py public-instances --heuristics all --timeout 20
\end{lstlisting}
Si l'algorithme 2 dépasse lui-même le délai, chaque heuristique est comparée
au statut attendu (fichier \texttt{expected.json} ou nom du fichier SATLIB).

\paragraph{Résultats.}
\begin{itemize}
  \item \textbf{SATLIB} (3200 instances uf20 à uf75 et uuf50 à uuf75 :
    2100 SAT, 1100 UNSAT) : les 16 heuristiques donnent le statut attendu
    sur toutes les instances, avec un modèle correct pour chaque instance
    SAT, aucun dépassement du délai de 10 s.
  \item \textbf{Instances publiques} (41 fichiers, délai de 20 s) : aucune
    erreur. Les 6 cas limites \texttt{edge\_*} et les tests d'interface
    passent pour toutes les heuristiques. L'algorithme 2 dépasse le délai
    sur public-03, public-11, public-12 et public-27 ; plusieurs heuristiques
    les résolvent (11 heuristiques sur 16 pour public-11, 12 pour
    public-27, toutes les variantes dynamiques pour public-12). Seule
    public-03, l'une des deux instances annoncées comme difficiles, n'est
    résolue par aucune. À l'inverse, quelques heuristiques dépassent le
    délai là où l'algorithme 2 conclut (H2 sur public-04 et public-21,
    H8s sur public-04, H3d, H5d, H7 et H8s sur public-06) : une heuristique n'est pas
    meilleure que le hasard sur toutes les instances.
\end{itemize}

\section{Premières mesures}
\label{sec:mesures}

L'analyse détaillée des performances est l'objet de la question 5. Pour
justifier nos choix, nous avons tout de même comparé les heuristiques sur
trois familles : 40 instances SATLIB uf50 (SAT), 40 instances uuf50 (UNSAT)
et 12 grilles Meowdoku générées de taille 8, 10 et 12 (délai de 10 s,
aucun dépassement). Le tableau donne la médiane du nombre de décisions et
la part du temps total passée dans le choix de branchement
(\texttt{Stats.heuristic\_time}).
\begin{center}
\begin{tabular}{lrrrrrr}
  \toprule
  & \multicolumn{2}{c}{uf50 (SAT)} & \multicolumn{2}{c}{uuf50 (UNSAT)}
  & \multicolumn{2}{c}{Meowdoku} \\
  Heuristique & décisions & \% heur. & décisions & \% heur.
  & décisions & \% heur. \\
  \midrule
  H0 (hasard)  & 160 & 10 & 432 & 9 & 14 & 7 \\
  H1           & 190 & 10 & 436 & 9 & 14 & 7 \\
  H2s / H2d    & 29 / 20 & 10 / 40 & 74 / 55 & 8 / 35 & 62 / 82 & 19 / 70 \\
  H3s / H3d    & 31 / 20 & 10 / 40 & 101 / 84 & 9 / 33 & 62 / 82 & 19 / 68 \\
  H4s / H4d    & 29 / \textbf{12} & 10 / 34 & 74 / \textbf{32} & 8 / 27 & 52 / 78 & 16 / 70 \\
  H5s / H5d    & 31 / 20 & 10 / 30 & 101 / 52 & 9 / 24 & 52 / 82 & 16 / 68 \\
  H6s / H6d    & 29 / 14 & 10 / 41 & 74 / 35 & 8 / 34 & 59 / 76 & 22 / 72 \\
  H7s / H7d    & 31 / 16 & 10 / 43 & 101 / 55 & 9 / 34 & 62 / 82 & 18 / 70 \\
  H8s / H8d    & 22 / 18 & 11 / 44 & 102 / 67 & 9 / 37 & \textbf{7} / 8 & 5 / 27 \\
  \bottomrule
\end{tabular}
\end{center}
\paragraph{Observations.}
\begin{itemize}
  \item Sur le 3-SAT aléatoire, toute heuristique fondée sur les
    occurrences divise le nombre de décisions par 5 à 13 par rapport au
    hasard, et le temps total par 5 à 11 (de 1,9 s à 0,17 s pour les 40
    instances uf50 avec H4d). H1 n'apporte rien : l'ordre des numéros de
    variables n'a pas de sens sur une formule aléatoire.
  \item Les variantes statiques H2s, H4s et H6s sont \emph{identiques} sur
    ces instances (de même que H3s, H5s et H7s) : toutes les clauses
    initiales ont 3 littéraux, donc les trois pondérations classent les
    variables de la même façon. Les différences n'apparaissent qu'avec les
    variantes dynamiques, quand les clauses raccourcissent.
  \item Le dynamique réduit encore le nombre de décisions (jusqu'à un
    facteur 2,5 pour H4 sur uuf50) mais 25 à 45\,\% du temps est alors
    passé à recalculer les scores, contre 10\,\% en statique. Sur ces petites
    instances le gain en temps reste faible ; sur l'instance publique
    public-12, en revanche, \emph{toutes} les variantes dynamiques
    concluent en moins de 20 s et \emph{aucune} variante statique.
    Le compromis \og coût de la mise à jour / pertinence \fg{} de l'énoncé
    sera étudié à la question 5.
  \item Sur Meowdoku, H2 à H7 font 4 à 6 fois \emph{plus} de décisions que
    le hasard, pour la raison exposée en section~\ref{sec:h8}. H8 divise
    par deux le nombre de décisions du hasard, tout en restant proche des
    meilleures heuristiques sur SATLIB.
\end{itemize}
Ces mesures confirment l'intérêt d'avoir implémenté toutes les variantes
plutôt que le minimum demandé : aucune heuristique n'est la meilleure
partout, et le classement dépend fortement de la structure des instances.

\end{document}
